\chapter{Introduction}

\paragraph{}
Plant diseases can have devastating effects on crops, leading to reduced overall yield \cite{kiranPlantLeafDisease2023, trivediAutomaticSegmentationPlant2022}. All approaches for developing disease-resistant crops or detecting plant pathogens require measurements of observable plant traits (phenotypes) \cite{gehanPlantCVV2Image2017}. Traditional plant phenotyping methods in plant pathogen detection require expertise. Moreover, these methods are destructive, labour-intensive, and time-consuming \cite{hussainRobustLeafDisease2024}. Therefore, traditional plant phenotyping methods are not suitable for monitoring large plant population sizes over time \cite{fahlgrenVersatilePhenotypingSystem2015, arnalbarbedoDigitalImageProcessing2013, zhangHighthroughputPhenotypingPlant2023, laxmanNoninvasiveQuantificationTomato2018}. As an alternative, non-destructive plant phenotyping methods using image processing have been developed \cite{furbankPhenomicsTechnologiesRelieve2011}, which are efficient, affordable, and allow larger samples of plants to be examined over time. However, these image-based solutions can be costly, making them inaccessible to less affluent plant growers \cite{lienLowcostOpensourcePlatform2019, tovarRaspberryPiPowered2018a, minerviniImageAnalysisNew2015}. Existing low-cost image-based phenotyping systems are patented, rendering the hardware and software of these systems inflexible and unmodifiable \cite{czedik-eysenbergPhenoBoxFlexibleAutomated2018}. Furthermore, many image processing tools in plant phenotyping are designed for specific plant species and environments and therefore cannot be generalized to other plant species and contexts \cite{arnalbarbedoDigitalImageProcessing2013, liReviewImagingTechniques2014}. Consequently, there is a need for low-cost, flexible, modifiable, generalizable, and extendable automated image processing tools for plant phenotyping in the field of plant pathogen detection.

\subparagraph{}
Recent developments in image-based plant phenotyping systems research show an increased development of plant phenotyping systems that use open-source hardware \cite{liDigitizationVisualizationGreenhouse2015, minerviniImageAnalysisNew2015} and software libraries \cite{fahlgrenVersatilePhenotypingSystem2015}. Despite these advances, several challenges still remain. In low-cost settings, commonly used strategies to manage the issues of large-scale image analysis and processing of phenotypic data include low-complexity traditional segmentation methods \cite{minerviniImageAnalysisNew2015}, the use of more affordable and straightforward imaging sensing technologies in low-cost off-the-shelf hardware \cite{tovarRaspberryPiPowered2018a}, and image compression techniques \cite{minerviniApplicationawareImageCompression2013}. However, each of these strategies comes with trade-offs in terms of precision, accuracy, and phenotypic analysis capabilities. Low-complexity traditional segmentation methods struggle to handle component-based phenotyping of specific organs in plants while requiring fewer computational resources \cite{choudhurySegmentationTechniquesChallenges2020}. RGB sensors are affordable and generate less raw data than other methods, such as hyperspectral imaging; however, their use is limited to the phenotyping of physiological plant traits \cite{liReviewImagingTechniques2014}. Image compression techniques can be used to reduce the amount of data transmission required for remote analysis but have a non-negligible impact on the accuracy of image segmentation \cite{minerviniApplicationawareImageCompression2013}. Therefore, a delicate balance must be struck between cost and numerous factors that affect the overall precision, accuracy, and measurability of phenotypic traits \cite{muller-linowPlantScreenMobile2019}.

\subparagraph{}
In response to the aforementioned challenges of image-based phenotyping solutions in plant pathogen detection, this paper aims to develop a binary support vector machine (SVM) \cite{vapnikNatureStatisticalLearning2000} classifier for the detection of tomato plant leaf diseases. The proposed classifier is trained using the publicly available PlantVillage dataset  \cite{hughesOpenAccessRepository2016}, and its overall reliability and performance are measured.

\subparagraph{}
The remainder of this paper is organized as follows: Section 2, the literature review, presents the fundamentals of plant phenotyping and its relation to plant pathogen detection. Particular emphasis is placed on existing image-based systems for plant pathogen detection. Thereafter, in section 3, the methodology, we discuss the proposed classification framework, detailing how images are processed, how the SVM classifier is trained, and how the overall system's performance and reliability are evaluated. In section 4, we discuss our experimental results after our SVM classifier has been trained and evaluated. Lastly, in section 5, we discuss possible extensions of our classifier, potential future research directions, and implications of our experimental results.